\documentclass[openany]{ctexbook}
\usepackage{amsmath,amssymb}
\usepackage[hidelinks]{hyperref}
\newcommand{\note}[1]{\footnote{译注：#1}}
\title{代数几何基础第五卷预备笔记\\推论10.4及其引导段}
\author{据用户上传英译编辑稿译出}
\date{2026年11月15日}
\begin{document}
\maketitle
\chapter*{范围与校订状态}
本批连续译出底本第40页“If in 10.2”起的未编号段，以及跨第40—41页的推论10.4全部文字，止于“quadratic”。下一接续处为第41页第11节“Grassmanians”标题及其首段。

本批仅对可核实的局部错误作有说明的校订；没有补写原稿未提供的证明。维数不等号原意不确，正文只采用下述光滑曲线范围内可核实的表述，原文及理由保留在脚注中。推论的整体存在性结论及所依赖的9.2尚未独立审定。此前8.8、9.5缺文等事项仍未解决。本批是独立增补，未合入原合订本。
\chapter*{推论10.4之前的未编号段}
若在10.2中取 $S=\operatorname{Spec}k$，并限于下述光滑曲线的情形，则所作的正则性假设等价于 $T=\varnothing$，因而 $Y_L\to X$ 是一个同构。\note{底本实际写作“$X$ is of dimension $\geq1$”。这一维数条件不能支持原结论：例如二维光滑仿射空间中由两个坐标给出的正则截面具有非空零点中心。对于下述光滑曲线，局部环是诺特环且维数至多1；非空中心处若存在长度2的正则列，则深度至少2，与深度不超过维数矛盾。反向采用10.2中允许空中心的正则截面（正则浸入）约定。故正文只写已核实的曲线情形，不把猜测的“$\leq1$”冒充已恢复的原稿。}因此，将此与9.2结合，得到：
\chapter*{推论10.4（10.2的推论）}
设 $X$ 是无限域 $k$ 上射影空间 $P$ 中的一条光滑、几何连通的曲线，且 $n\geq2$。则存在 $P$ 上的一个 $n$ 次形式的线性铅笔，它定义一个态射
\[
X\longrightarrow\mathbf P^1,
\]
并具有下述性质：这个态射一般是 étale 的，次数为 $d$；对于 $\mathbf P^1$ 的每一个几何点 $s$，$X_s$ 或者在代数闭域 $k'=k(s)$ 上是 étale 的，或者在 $k'$ 上同构于 $d-2$ 个概形 $\operatorname{Spec}k'$ 与概形
\[
I_{k'}=\operatorname{Spec}\bigl(k'[t]/(t^2)\bigr)
\]
的不交和。\note{底本在“$d-2$ 个概形”处印作 $\operatorname{Spec}k$；这里描述的是 $k'$ 上的几何纤维，故按其底域更正为 $\operatorname{Spec}k'$。$k'=k(s)$ 沿用底本的几何点记号；这里的 $s$ 是取值于代数闭域的几何点，不能误作底空间任意点的未代数闭剩余域。}

用前辈们的语言来说：每个这样的几何纤维中至多有一个分歧点，而且它是“二次的”。\note{底本末句为“there is at most one point of ramification and it is ‘quadratic.’”这里“至多一个”承接前句，逐个几何纤维而言，不是说整个态射总共只有一个分歧点。“二次的”按上面的局部纤维 $\operatorname{Spec}(k'[t]/(t^2))$ 理解；不额外断言该分歧在任意特征下都是温和分歧。底本未在此另给证明，本译文也未独立验证铅笔的存在性及推论的全部假设。}
\end{document}
